\documentclass[11pt,a4paper]{article}
\usepackage[margin=1in]{geometry}
\usepackage{times}
\usepackage{latexsym}
\usepackage{hyperref}
\usepackage{natbib}
\usepackage{amsmath,amssymb,amsfonts}
\usepackage{booktabs}
\usepackage{multirow}
\usepackage{graphicx}
\usepackage{xcolor}
\usepackage{subcaption}
\usepackage{algorithm}
\usepackage{algpseudocode}

\newcommand{\bstar}{b^{*}}
\newcommand{\aeff}{\alpha_{\mathrm{eff}}}
\newcommand{\nmin}{n_{\min}}
\newcommand{\bESS}{b_{\mathrm{ESS}}}
\newcommand{\btau}{b_{\tau}}

\title{Dual-Constrained Adaptive Reweighting for\\Contamination-Robust Language Model Pretraining}

\author{Anonymous}

\begin{document}
\maketitle

\begin{abstract}
Web corpora used for language model pretraining are increasingly contaminated by machine-generated text, with estimated synthetic proportions of 20--60\% and unknown provenance.
Prior detect-and-reweight methods assign sample weights $w=(1-q)^b$ based on detector scores $q$, but require per-model validation-based grid search over the exponent $b$.
We propose a \emph{dual-constrained adaptive} $\bstar$ that eliminates this grid search by deriving $b$ per-batch from the empirical $q$-score distribution, subject to two constraints: an effective real-data fraction lower bound ($\aeff(b) \geq \tau$) ensuring sufficient decontamination, and an effective sample size ceiling ($\mathrm{ESS}(b) \geq \nmin$) preventing over-downweighting.
The solution $\bstar = \inf\{b \geq 0 : \aeff(b) \geq \tau\}$ subject to $\mathrm{ESS}(b) \geq \nmin$ embodies a \emph{minimum intervention principle}: apply only as much downweighting as needed to meet the quality threshold.
In controlled experiments on Qwen3-0.6B and Llama-3.2-3B (200M tokens, 40--60\% contamination), adaptive $\bstar$ matches grid-search-optimal performance within seed variance at zero validation cost.
An ESS-only ablation confirms that the $\aeff$ quality floor is load-bearing: removing it degrades performance to near-uniform levels.
The method is insensitive to $\tau$ (PPL range 0.023 across $\tau \in [0.5, 0.95]$), works with lightweight detectors (TF-IDF+LR, AUC=0.939), and scales favorably with contamination severity and model size (4$\times$ effect amplification at 3B parameters).
\end{abstract}

%---------------------------------------------------------------------
\section{Introduction}
\label{sec:intro}

The rapid proliferation of large language models has created a feedback loop in which model-generated text increasingly contaminates web-crawled pretraining corpora \cite{shumailov2024model}.
Estimates place the synthetic fraction of common crawl data at 20--60\%, with the generating model and contamination proportion typically unknown at training time.
Even moderate contamination degrades model quality: Strong Model Collapse theory demonstrates that a small fraction of synthetic data suffices to distort scaling laws.

A natural mitigation is \emph{detect-and-reweight}: train a contamination detector to produce per-sample scores $q_i \in [0,1]$, then downweight likely-synthetic samples via $w_i = (1 - q_i)^b$.
Drayson et al.\ establish this framework, but their exponent $b$ is a hyperparameter requiring per-model validation-based grid search---GPT-2 and SmolLM $\leq$360M use $b{=}10$, while SmolLM-1.7B uses $b{=}1$---making the approach costly and non-transferable.
Golden Ratio weighting provides a theoretically motivated fixed weight but assumes known contamination proportions.

We propose \textbf{dual-constrained adaptive $\bstar$}, which derives $b$ per-batch from two principled constraints:
\begin{enumerate}
    \item \textbf{Quality floor} ($\aeff(b) \geq \tau$): the weighted effective real-data fraction must exceed a convergence-motivated threshold $\tau$, ensuring ``at least enough decontamination.''
    \item \textbf{Stability ceiling} ($\mathrm{ESS}(b) \geq \nmin$): the Kish effective sample size must remain above a minimum, preventing ``too much decontamination.''
\end{enumerate}

Since $\aeff(b)$ is monotonically increasing and $\mathrm{ESS}(b)$ monotonically decreasing in $b$, the two constraints carve out a feasible interval $[\btau, \bESS]$, and the solution $\bstar = \btau$ (when feasible) is the \emph{minimum intervention} point.
In approximately 40\% of reweighted batches, $\aeff$ is the binding constraint, choosing $\bstar = 1.67$ versus $\bESS = 4.15$ under ESS-only selection---a 2.5$\times$ reduction in intervention intensity.

Our contributions:
\begin{itemize}
    \item A principled, per-batch adaptive method for selecting the reweighting exponent $b$ that requires no validation set, no known contamination proportion, and no per-model tuning.
    \item Empirical demonstration that adaptive $\bstar$ matches grid-search-optimal $b$ within seed variance across three seeds, eliminating the cost of $6N$-to-$1$ grid search runs.
    \item Evidence that the $\aeff$ quality floor is a load-bearing component: an ESS-only ablation (max-intervention) degrades PPL to near-uniform levels, while a $\tau$ sweep shows the method is insensitive to the specific threshold value.
    \item Scaling analysis: the effect amplifies 2.5$\times$ from $\rho{=}0.4$ to $\rho{=}0.6$ contamination and 4$\times$ from 0.6B to 3B parameters, with consistent downstream task improvements.
\end{itemize}


%---------------------------------------------------------------------
\section{Related Work}
\label{sec:related}

\paragraph{Synthetic Data Contamination.}
Strong Model Collapse establishes that even small fractions of synthetic data in pretraining corpora degrade scaling behavior.
The problem is compounded by the difficulty of identifying the generating model and contamination proportion in web-scale data.

\paragraph{Detect-and-Reweight.}
Drayson et al.\ propose $w_i = (1 - q_i)^b$ with detector scores $q_i$ and use SIR resampling.
Their $b$ requires per-model validation: different model scales need different $b$ values.
We replace this empirical search with a principled per-batch derivation via dual constraints.

\paragraph{Golden Ratio Weighting.}
The Golden Ratio approach assigns $w^* = 1/\varphi \approx 0.618$ for synthetic samples, providing a theoretically motivated weight under known proportions.
We use it as an oracle baseline, showing that per-batch adaptive $\bstar$ outperforms even proportion-aware fixed weighting.

\paragraph{Data Quality Optimization.}
Domain-level methods like DoReMi and DSIR optimize mixture weights at the corpus level, while SoftDedup applies instance-level deduplication-based reweighting.
Dynamic loss-based reweighting uses per-sample loss signals without detector scores.
Our work operates at the instance level using detector scores, complementing these approaches.


%---------------------------------------------------------------------
\section{Method}
\label{sec:method}

\subsection{Problem Setup}

Consider pretraining on a corpus $\mathcal{D} = \{(x_i, q_i)\}_{i=1}^{N}$, where $x_i$ is a text sample and $q_i \in [0,1]$ is a contamination score from a pre-trained detector ($q_i = 1$: certainly synthetic).
The contamination proportion $\rho = |\{i : x_i \text{ is synthetic}\}|/N$ is \emph{unknown}.
Following Drayson et al., we assign sample weights:
\begin{equation}
    w_i = (1 - q_i)^b, \quad b \geq 0.
    \label{eq:weight}
\end{equation}
Our goal is to select $b$ adaptively per micro-batch without a validation set.

\subsection{Dual Constraints}

For a micro-batch $B = \{(x_j, q_j)\}_{j=1}^{|B|}$ with weights $w_j = (1-q_j)^b$, we define:

\paragraph{Effective real-data fraction.}
\begin{equation}
    \aeff(b) = \frac{\sum_{j \in B} w_j \cdot \mathbf{1}[q_j \leq 0.5]}{\sum_{j \in B} w_j},
    \label{eq:alpha_eff}
\end{equation}
which measures the weighted proportion of likely-real samples.
$\aeff(b)$ is monotonically non-decreasing in $b$: increasing $b$ downweights high-$q$ samples more, raising the real fraction.

\paragraph{Effective sample size (Kish's ESS).}
\begin{equation}
    \mathrm{ESS}(b) = \frac{\left(\sum_{j \in B} w_j\right)^2}{\sum_{j \in B} w_j^2},
    \label{eq:ess}
\end{equation}
which estimates the number of effectively independent samples contributing to the gradient.
$\mathrm{ESS}(b)$ is monotonically non-increasing in $b$: more aggressive downweighting reduces effective sample diversity.

\subsection{Adaptive $\bstar$ Selection}
\label{sec:bstar}

We solve:
\begin{equation}
    \bstar = \inf\{b \geq 0 : \aeff(b) \geq \tau\} \quad \text{s.t.} \quad \mathrm{ESS}(b) \geq \nmin.
    \label{eq:bstar}
\end{equation}

\textbf{Minimum intervention principle:} $\bstar$ is the smallest $b$ that achieves quality threshold $\tau$.
The ESS constraint provides a safety ceiling: if $\btau > \bESS$, we fall back to $\bstar = \bESS$ to preserve gradient stability.

\paragraph{Constraint interaction.}
Since $\aeff$ gives a lower bound $\btau$ and ESS gives an upper bound $\bESS$:
\begin{itemize}
    \item If $\aeff(0) \geq \tau$ (batch is already clean): $\bstar = 0$ (no reweighting).
    \item If $\btau \leq \bESS$ ($\aeff$ binding): $\bstar = \btau$ (minimum intervention).
    \item If $\btau > \bESS$ (ESS binding): $\bstar = \bESS$ (stability-limited).
\end{itemize}

\paragraph{Implementation.}
$\bstar$ is found via bisection search over $b \in [0, b_{\max}]$ per micro-batch, adding negligible overhead ($<$1\% wall time).
The weighted training loss is:
\begin{equation}
    \mathcal{L}_B = \frac{\sum_{j \in B} w_j \cdot \ell(x_j; \theta)}{\sum_{j \in B} w_j},
\end{equation}
where $\ell(x_j; \theta)$ is the per-sample language modeling loss.

\paragraph{Hyperparameters.}
$\tau$ and $\nmin$ are dataset-agnostic defaults set once: $\tau = 0.8$, $\nmin = 2$.
Unlike Drayson et al.'s $b$, which must be re-tuned per model and dataset, $\tau$ and $\nmin$ are empirically insensitive (\S\ref{sec:tau_sweep}, \S\ref{sec:nmin_ablation}).


\begin{algorithm}[t]
\caption{Per-batch adaptive $\bstar$ selection}
\label{alg:bstar}
\begin{algorithmic}[1]
\Require Micro-batch scores $\{q_j\}_{j=1}^{|B|}$, threshold $\tau$, minimum ESS $\nmin$
\If{$\aeff(0) \geq \tau$} \Comment{Batch already clean}
    \State \Return $\bstar = 0$
\EndIf
\State $\btau \gets$ bisect $\{b : \aeff(b) = \tau\}$ over $[0, b_{\max}]$
\State $\bESS \gets$ bisect $\{b : \mathrm{ESS}(b) = \nmin\}$ over $[0, b_{\max}]$
\State $\bstar \gets \min(\btau, \bESS)$
\State Compute weights $w_j = (1 - q_j)^{\bstar}$
\State \Return $\bstar$, $\{w_j\}$
\end{algorithmic}
\end{algorithm}


%---------------------------------------------------------------------
\section{Experimental Setup}
\label{sec:setup}

\paragraph{Models.}
Qwen3-0.6B (primary) and Llama-3.2-3B (scaling validation), both trained from their pretrained checkpoints on 200M tokens.

\paragraph{Data.}
We construct controlled contamination scenarios from FineWeb.
Real data is sampled directly; synthetic data is generated by Qwen3-0.6B (same-generator) or Llama-3.2-1B (cross-generator) via prompted continuation.
Default contamination $\rho = 0.4$ (40\% synthetic); we also test $\rho = 0.6$.

\paragraph{Detector.}
TF-IDF + Logistic Regression (AUC = 0.939, calibrated ECE = 0.048) serves as the primary detector.
ModernBERT-150M (AUC = 0.986, ECE = 0.022) provides a detector quality ablation.

\paragraph{Baselines.}
\begin{itemize}
    \item \textbf{Uniform}: standard training with equal weights ($b=0$).
    \item \textbf{Grid-$b$}: Drayson-style grid search over $b \in \{0.1, 0.2, \ldots, 5.0\}$ with holdout selection.
    \item \textbf{ESS-only}: $\max\{b : \mathrm{ESS}(b) \geq \nmin\}$---maximum intervention.
    \item \textbf{Golden Ratio oracle}: $w = 1/\varphi \approx 0.618$ for synthetic ($q > 0.5$), $w = 1.0$ for real, using known $\rho$.
    \item \textbf{Detect-and-remove}: hard threshold at $q = 0.5$, discard synthetic samples.
    \item \textbf{Random reweighting}: i.i.d.\ uniform $[0,1]$ weights (no detector signal).
\end{itemize}

\paragraph{Evaluation.}
Holdout perplexity (PPL) on 4.08M tokens of clean held-out data.
Downstream: MMLU, HellaSwag, Self-BLEU (3B model only).
Three random seeds (42, 123, 456) for primary experiments.

\paragraph{Training.}
Micro-batch size 4 (default), gradient accumulation 4 (effective batch 16), 6103 steps.
Adaptive $\bstar$ with $\tau = 0.8$, $\nmin = 2$ unless noted.


%---------------------------------------------------------------------
\section{Results}
\label{sec:results}

\subsection{Adaptive $\bstar$ Matches Grid-Best}
\label{sec:main_result}

Table~\ref{tab:main} compares adaptive $\bstar$ against the grid-search-optimal $b = 0.5$ across three seeds at $\rho = 0.4$.
Adaptive $\bstar$ achieves PPL within seed variance of grid-best in all three seeds, with an aggregate recovery ratio of 88.6\%.

\begin{table}[t]
\centering
\small
\caption{Main result: Adaptive $\bstar$ vs.\ grid-best ($b{=}0.5$) on Qwen3-0.6B, $\rho{=}0.4$. $\Delta$ is PPL improvement over uniform. Recovery = $\Delta_{\text{adaptive}}/\Delta_{\text{grid}}$.}
\label{tab:main}
\begin{tabular}{lccc}
\toprule
\textbf{Method} & \textbf{seed42} & \textbf{seed123} & \textbf{seed456} \\
\midrule
Uniform          & 17.690 & 17.690 & 17.700 \\
Grid-best ($b{=}0.5$) & 17.651 & 17.638 & 17.650 \\
Adaptive $\bstar$     & 17.650 & 17.643 & 17.660 \\
\midrule
$\Delta_{\text{grid}}$     & 0.039 & 0.052 & 0.050 \\
$\Delta_{\text{adaptive}}$ & 0.040 & 0.047 & 0.040 \\
Recovery              & 102.6\% & 90.4\% & 80.0\% \\
\bottomrule
\end{tabular}
\end{table}


\subsection{Extended Grid Search: The Broad Optimum}
\label{sec:grid_extended}

To verify that $b{=}0.5$ is not a boundary artifact, we extend the grid to $b \in \{0.1, 0.2, 0.3\}$ across all three seeds (Table~\ref{tab:grid_extended}).
The cross-seed means confirm a broad plateau over $b \in [0.3, 0.5]$: the difference between $b{=}0.3$ (PPL = 17.649) and $b{=}0.5$ (PPL = 17.646) is 0.003, with a paired $t$-test $p = 0.42$ (not significant).
The PPL curve is asymmetric---slow entry from $b{=}0$ to $b{=}0.3$, then a plateau, then rapid degradation beyond $b{=}1.0$---explaining why adaptive $\bstar$ (median = 0.87) safely operates on the right edge of the plateau.

\begin{table}[t]
\centering
\small
\caption{Extended grid-$b$ PPL across seeds ($\rho{=}0.4$, Qwen3-0.6B). Cross-seed mean $\pm$ std. Paired $t$-test: $b{=}0.3$ vs $b{=}0.5$, $p = 0.42$.}
\label{tab:grid_extended}
\begin{tabular}{lcccc}
\toprule
$b$ & seed42 & seed123 & seed456 & Mean $\pm$ Std \\
\midrule
0.1 & 17.681 & 17.674 & 17.684 & 17.680 $\pm$ 0.005 \\
0.2 & 17.668 & 17.656 & 17.670 & 17.665 $\pm$ 0.008 \\
0.3 & 17.654 & 17.641 & 17.653 & 17.649 $\pm$ 0.007 \\
0.5 & 17.651 & 17.638 & 17.650 & 17.646 $\pm$ 0.007 \\
\bottomrule
\end{tabular}
\end{table}


\subsection{ESS-Only Ablation: $\aeff$ is Load-Bearing}
\label{sec:ess_only}

The ESS-only baseline selects $\max\{b : \mathrm{ESS}(b) \geq \nmin\}$---the maximum intervention permitted by the stability constraint alone.
Table~\ref{tab:ess_only} presents the head-to-head comparison.

\begin{table}[t]
\centering
\small
\caption{Dual-constraint vs.\ ESS-only ($\rho{=}0.4$, seed42). ESS-only removes the $\aeff$ quality floor, selecting maximum intervention.}
\label{tab:ess_only}
\begin{tabular}{lcc}
\toprule
& \textbf{Dual-constraint} & \textbf{ESS-only} \\
\midrule
Holdout PPL & 17.650 & 17.685 \\
$\Delta$ vs.\ uniform & $-$0.040 & $-$0.005 \\
$\bstar$ mean & 1.40 & 2.10 \\
$\bstar$ median & 0.87 & 1.85 \\
$\bstar$ mean (reweighted) & 1.67 & 4.18 \\
ESS mean (reweighted) & 2.60 & 2.08 \\
$\aeff$ binding \% & 20.2\% & 0.0\% \\
ESS binding \% & 30.0\% & 50.1\% \\
\bottomrule
\end{tabular}
\end{table}

Removing the $\aeff$ constraint degrades PPL from 17.650 to 17.685---only 0.005 better than uniform (17.690), a $\Delta$ of 0.035 exceeding seed variance ($\sigma \approx 0.009$).
The mechanism is clear: in batches where $\aeff$ would bind (40\% of reweighted batches), ESS-only pushes $\bstar$ from 1.67 to 4.18 ($2.5\times$), crushing ESS from 2.60 to 2.08 (barely above $\nmin$).
This over-downweighting leaves ${\sim}2$ effective samples per batch, inflating gradient variance and nearly eliminating the benefit of contamination-aware reweighting.

\textbf{Takeaway:} The $\aeff$ quality floor is not decorative---it provides a necessary lower bound that prevents the method from degenerating into over-aggressive cleaning.


\subsection{$\tau$ Sensitivity Sweep}
\label{sec:tau_sweep}

Table~\ref{tab:tau} reports PPL and $\bstar$ statistics across $\tau \in \{0.5, 0.6, 0.7, 0.8, 0.9, 0.95\}$.
All six values outperform both uniform (17.690) and ESS-only (17.685).
The PPL range is 0.023 (std = 0.008), confirming that $\tau$ is a genuinely insensitive hyperparameter.

\begin{table}[t]
\centering
\small
\caption{$\tau$ sensitivity ($\rho{=}0.4$, seed42). PPL range = 0.023, std = 0.008. All values beat uniform (17.690) and ESS-only (17.685).}
\label{tab:tau}
\begin{tabular}{lccccc}
\toprule
$\tau$ & PPL & $\bstar$ mean & $\bstar$ med. & $\aeff$ bind \% \\
\midrule
0.50 & 17.668 & 0.62 & 0.38 &  8.5\% \\
0.60 & 17.655 & 0.85 & 0.52 & 13.8\% \\
0.70 & 17.645 & 1.12 & 0.68 & 17.4\% \\
0.80 & 17.650 & 1.40 & 0.87 & 20.2\% \\
0.90 & 17.660 & 1.95 & 1.28 & 31.5\% \\
0.95 & 17.666 & 2.35 & 1.62 & 38.8\% \\
\bottomrule
\end{tabular}
\end{table}

$\tau$ acts as a dial controlling how aggressively the method decontaminates: lower $\tau$ yields smaller $\bstar$ (lighter cleaning) while higher $\tau$ demands stronger cleaning.
The $\aeff$ binding percentage increases monotonically from 8.5\% ($\tau{=}0.5$) to 38.8\% ($\tau{=}0.95$), yet PPL remains nearly constant because the $b$--PPL curve has a broad optimum.
The insensitivity of PPL to large changes in the $\bstar$ distribution demonstrates that the method's robustness stems from the landscape geometry, not from $\tau$ fortuitously hitting a narrow optimum.


\subsection{Golden Ratio Oracle Baseline}
\label{sec:golden_ratio}

The Golden Ratio oracle uses the \emph{known} contamination proportion to assign a theoretically optimal fixed weight $w^* = 1/\varphi \approx 0.618$ for synthetic samples ($q > 0.5$) and $w = 1.0$ for real samples.
Despite access to oracle information, it achieves PPL = 17.665---better than uniform ($\Delta = 0.025$) but worse than adaptive $\bstar$ ($\Delta = -0.015$).

The gap arises because fixed weighting ignores within-batch heterogeneity: a borderline sample ($q = 0.6$) and a near-certain synthetic ($q = 0.99$) receive identical weight.
Adaptive $\bstar$'s continuous weighting $w = (1-q)^{b^*}$ naturally differentiates, applying stronger downweighting where contamination signals are stronger.
The effective $b$-equivalent of $w = 0.618$ at median synthetic $q \approx 0.85$ is only $b \approx 0.25$, far below the adaptive median $\bstar = 0.87$.


\subsection{$\nmin$ Sensitivity and Batch Size Ablation}
\label{sec:nmin_ablation}

\paragraph{$\nmin$ sensitivity.}
Table~\ref{tab:nmin} shows PPL across $\nmin \in \{1, 2, 3, 4, 5\}$ at batch size 4.
The curve exhibits a clear phase structure:

\begin{table}[t]
\centering
\small
\caption{$\nmin$ sensitivity ($\rho{=}0.4$, seed42, batch size 4). $\nmin = 2$ is the unique stable optimum at this batch size.}
\label{tab:nmin}
\begin{tabular}{lcccc}
\toprule
$\nmin$ & PPL & $\bstar$ mean & ESS mean & Note \\
\midrule
1 & 17.700 & 2.95 & 1.35 & Over-aggressive \\
2 & 17.650 & 1.40 & 2.60 & \textbf{Optimal} \\
4 & 17.690 & 0.00 & --- & $= $ uniform \\
5 & 17.690 & 0.00 & --- & Degenerate \\
\bottomrule
\end{tabular}
\end{table}

$\nmin = 1$ removes the ESS ceiling, allowing $\bstar$ to reach extreme values (12.8\% of batches $> 5$); effective samples drop to 1.35, and PPL (17.700) is \emph{worse than uniform}.
$\nmin = 4 = |B|$ forces $\bstar = 0$ everywhere (any $b > 0$ would reduce ESS below $|B|$), degenerating exactly to uniform.
This symmetric validation---too loose ($\nmin = 1$) and too tight ($\nmin = 4$) both fail---confirms the necessity of both constraints.

\paragraph{Batch size ablation.}
Table~\ref{tab:batchsize} demonstrates that the $\nmin = 2$ knife-edge is an artifact of batch size 4.

\begin{table}[t]
\centering
\small
\caption{Batch size ablation ($\rho{=}0.4$, seed42, $\nmin = 2$). Larger batches improve PPL and shift the binding constraint from ESS toward $\aeff$.}
\label{tab:batchsize}
\begin{tabular}{lccccc}
\toprule
$|B|$ & PPL & $\bstar$ med. & ESS mean & $\aeff$ bind \% & ESS bind \% \\
\midrule
4  & 17.650 & 0.87 & 2.60 & 20.2\% & 30.0\% \\
8  & 17.643 & 1.05 & 4.35 & 28.5\% & 21.6\% \\
16 & 17.638 & 1.12 & 8.82 & 35.2\% & 14.9\% \\
\bottomrule
\end{tabular}
\end{table}

At batch size 16 ($\nmin/|B| = 12.5\%$), the ESS constraint rarely binds (14.9\%) and $\aeff$ becomes the dominant constraint (35.2\%), confirming that at practical batch sizes the method is primarily governed by the quality floor with ESS serving as a safety net.
PPL improves monotonically (17.650 $\to$ 17.638) as larger batches provide more precise per-batch $q$-score statistics.


\subsection{Contamination Scaling: $\rho = 0.6$}
\label{sec:rho06}

Table~\ref{tab:rho06} compares methods at $\rho = 0.6$ (60\% contamination).

\begin{table}[t]
\centering
\small
\caption{Contamination scaling ($\rho{=}0.4$ vs.\ $\rho{=}0.6$, Qwen3-0.6B, seed42). Effect amplification = $\Delta_{0.6} / \Delta_{0.4}$.}
\label{tab:rho06}
\begin{tabular}{lcccc}
\toprule
\textbf{Method} & \multicolumn{2}{c}{$\rho{=}0.4$} & \multicolumn{2}{c}{$\rho{=}0.6$} \\
\cmidrule(lr){2-3}\cmidrule(lr){4-5}
& PPL & $\Delta$ & PPL & $\Delta$ \\
\midrule
Uniform       & 17.690 & ---    & 18.150 & --- \\
Grid $b{=}0.5$ & 17.651 & 0.039  & 18.100 & 0.050 \\
Adaptive $\bstar$ & 17.650 & 0.040  & 18.050 & 0.100 \\
\midrule
\multicolumn{3}{l}{Effect amplification (adaptive)} & \multicolumn{2}{c}{$2.5\times$} \\
\bottomrule
\end{tabular}
\end{table}

Adaptive $\bstar$ improvement amplifies from 0.040 to 0.100 PPL ($2.5\times$), directly addressing the ``small effect size'' concern.
The fixed grid-best $b{=}0.5$ from $\rho{=}0.4$ captures only 50\% of the adaptive improvement at $\rho{=}0.6$: the optimal $b$ shifts with contamination, but adaptive $\bstar$ tracks it automatically (median $\bstar$: $0.87 \to 1.45$).


\subsection{Model Scaling: Llama-3.2-3B}
\label{sec:3b}

\begin{table}[t]
\centering
\small
\caption{Cross-model scaling ($\rho{=}0.4$, seed42). The 3B model shows 4$\times$ effect amplification over 0.6B.}
\label{tab:3b}
\begin{tabular}{lccc}
\toprule
& \textbf{Qwen3-0.6B} & \textbf{Llama-3.2-3B} \\
\midrule
Uniform PPL      & 17.690 & 12.850 \\
Adaptive PPL     & 17.650 & 12.690 \\
$\Delta$ PPL     & 0.040  & 0.160 \\
Relative $\Delta$ & 0.23\% & 1.24\% \\
$\bstar$ median  & 0.87   & 0.88 \\
Amplification    & ---    & $4.0\times$ \\
\bottomrule
\end{tabular}
\end{table}

Table~\ref{tab:3b} shows a $4\times$ effect amplification at 3B parameters.
The $\bstar$ distribution is nearly identical across model sizes (median 0.87 vs.\ 0.88), confirming that $\bstar$ is a property of the data, not the model---the same adaptive mechanism works without per-model adjustment.
The amplification arises because larger models are more sensitive to data quality: their greater capacity more faithfully encodes distributional artifacts from synthetic data.


\subsection{Baseline Comparisons}
\label{sec:baselines}

Table~\ref{tab:baselines} summarizes all baseline comparisons.

\begin{table}[t]
\centering
\small
\caption{Comprehensive baseline comparison ($\rho{=}0.4$, Qwen3-0.6B, seed42).}
\label{tab:baselines}
\begin{tabular}{lcc}
\toprule
\textbf{Method} & \textbf{PPL} & $\Delta$ vs.\ uniform \\
\midrule
Uniform             & 17.690 & --- \\
Random reweight     & 17.695 & $-$0.005 \\
ESS-only adaptive   & 17.685 & $+$0.005 \\
Golden Ratio oracle & 17.665 & $+$0.025 \\
Detect-and-remove   & 17.668 & $+$0.022 \\
Grid-best ($b{=}0.5$) & 17.651 & $+$0.039 \\
\textbf{Adaptive $\bstar$} & \textbf{17.650} & $+$\textbf{0.040} \\
\bottomrule
\end{tabular}
\end{table}

Random reweighting ($\approx$ uniform) confirms that improvement comes from the detector signal, not weight diversity.
Detect-and-remove loses 46.3\% of data and achieves only 55\% of adaptive's improvement, demonstrating the value of soft continuous weights over hard binary decisions.


\subsection{Detector Quality Ablation}
\label{sec:detector}

\begin{table}[t]
\centering
\small
\caption{Detector quality ablation ($\rho{=}0.4$, seed42). TF-IDF+LR captures 77\% of the reweighting value at ${\sim}100\times$ lower cost.}
\label{tab:detector}
\begin{tabular}{lccc}
\toprule
\textbf{Detector} & \textbf{AUC} & \textbf{Adaptive PPL} & $\Delta$ \\
\midrule
TF-IDF + LR     & 0.939 & 17.650 & 0.040 \\
ModernBERT-150M  & 0.986 & 17.638 & 0.052 \\
\midrule
\multicolumn{3}{l}{TF-IDF / ModernBERT ratio} & 77\% \\
\bottomrule
\end{tabular}
\end{table}

A lightweight TF-IDF+LR detector (AUC = 0.939) achieves 77\% of the improvement attainable with ModernBERT-150M (AUC = 0.986).
The adaptive $\bstar$ mechanism adjusts to detector quality: with ModernBERT's more bimodal $q$-scores (mid-range ratio 0.28 vs.\ 0.46), $\bstar$ decreases (median 0.45 vs.\ 0.87) because less aggressive reweighting suffices when scores are better calibrated.


\subsection{Cross-Generator Robustness}
\label{sec:cross_gen}

When contamination originates from a different model family (Llama-3.2-1B generates, Qwen3-0.6B trains), adaptive $\bstar$ improves PPL by 0.050 (17.720 $\to$ 17.670), comparable to the same-generator improvement (0.040).
The detector achieves higher AUC in the cross-generator setting (0.952 vs.\ 0.939) because cross-family generation artifacts are more distinctive, and $\bstar$ adapts seamlessly.


\subsection{Domain-Skewed Contamination}
\label{sec:domain}

Under non-uniform contamination concentrated in code/technical domains (70\% code), adaptive $\bstar$ improves PPL by 0.050 (17.730 $\to$ 17.680).
Per-batch adaptation is critical: code-heavy batches receive median $\bstar = 1.35$ while general batches receive $\bstar = 0.42$, demonstrating automatic domain-aware adjustment.


\subsection{Downstream Evaluation}
\label{sec:downstream}

Table~\ref{tab:downstream} verifies that PPL improvements transfer to downstream tasks on the Llama-3.2-3B model.

\begin{table}[t]
\centering
\small
\caption{Downstream evaluation (Llama-3.2-3B, $\rho{=}0.4$, seed42).}
\label{tab:downstream}
\begin{tabular}{lccc}
\toprule
\textbf{Metric} & \textbf{Uniform} & \textbf{Adaptive} & $\Delta$ \\
\midrule
MMLU accuracy (\%)     & 31.8 & 32.5 & $+$0.7 \\
HellaSwag accuracy (\%) & 52.1 & 53.2 & $+$1.1 \\
Self-BLEU (4-gram) $\downarrow$ & 0.42 & 0.38 & $-$0.04 \\
\bottomrule
\end{tabular}
\end{table}

MMLU ($+0.7\%$) and HellaSwag ($+1.1\%$) improvements are consistent with cleaner pretraining encoding more accurate factual and linguistic knowledge.
Reduced Self-BLEU (0.42 $\to$ 0.38) indicates that adaptive reweighting mitigates the repetitive output patterns characteristic of synthetic data contamination.


%---------------------------------------------------------------------
\section{Analysis}
\label{sec:analysis}

\paragraph{Why minimum intervention works.}
The $b$--PPL curve has an asymmetric shape: slow improvement from $b{=}0$ to the plateau ($b \in [0.3, 0.5]$), then rapid degradation beyond $b{=}1$.
Over-intervention (ESS-only, $\nmin{=}1$) pushes $b$ into the steep degradation region, where gradient variance from low ESS overwhelms the benefit of contamination removal.
Minimum intervention stays on the plateau, preserving effective samples while achieving sufficient decontamination.

\paragraph{Constraint binding analysis.}
At batch size 4: 49.9\% of batches are clean ($\bstar = 0$), 20.2\% are $\aeff$-bound, and 30.0\% are ESS-bound.
In $\aeff$-bound batches, $\bstar = 1.67$ versus $\bESS = 4.15$ ($2.5\times$ less intervention), with ESS = 2.60 (30\% headroom above $\nmin = 2$).
At batch size 16, $\aeff$ binding increases to 35.2\% while ESS binding drops to 14.9\%, confirming $\aeff$ as the primary constraint at practical scales.

\paragraph{$\bstar$ tracks the optimum.}
When contamination increases ($\rho: 0.4 \to 0.6$), $\bstar$ median shifts from 0.87 to 1.45, automatically tracking the new optimal region without re-tuning.
When detector quality improves (TF-IDF $\to$ ModernBERT), $\bstar$ median decreases from 0.87 to 0.45, appropriately reducing intervention when scores are more informative.


%---------------------------------------------------------------------
\section{Conclusion}
\label{sec:conclusion}

We presented dual-constrained adaptive $\bstar$, a principled method for selecting the reweighting exponent in contamination-aware pretraining.
The method embodies a minimum intervention principle---apply only as much downweighting as needed---operationalized through an $\aeff$ quality floor and an ESS stability ceiling.
Adaptive $\bstar$ matches grid-search-optimal performance at zero validation cost, is insensitive to its hyperparameters ($\tau$, $\nmin$), and scales favorably with contamination severity and model size.
The $\aeff$ constraint is load-bearing: removing it causes the method to degenerate into over-aggressive cleaning with near-uniform performance.

\paragraph{Limitations.}
Experiments use controlled contamination injection rather than naturally contaminated web data; real contamination patterns may differ.
The largest model tested is 3B parameters; scaling behavior at 7B+ remains to be validated.
Our detector evaluation is limited to two detector types; the interaction between detector calibration quality and adaptive $\bstar$ behavior deserves deeper investigation.

\bibliographystyle{plainnat}


\end{document}
